%************************************************
\chapter{Introduction}\label{ch:introduction}
%************************************************

\note{Definir proceso estocástico X con trayectorias en L2 y quizás después decir que suponemos que son continuas. Observado en un grid denso.}

\note{Contributions? chapter-by-chapter?}

\note{Notación común para producto escalar L2 \(\dotprod{\cdot}{\cdot}_2\). Quizás en el glosario de símbolos ponerlo junto a las otras posibles notaciones (norma de matrices,..)}

The problem of predicting a scalar response from a functional covariate is one that has gained traction over the last few decades, as more and more real-world data is being generated with an ever-increasing level of granularity in the measurements. While in principle the functional data could simply be regarded as a discretized vector in a very high dimension, there are often many advantages in taking into account its functional nature, ranging from modeling the correlation among points that are close in the domain, to extracting information that may be hidden in the derivatives of the function in question. As a consequence, numerous proposals have arisen on how to suitably deal with functional data, all of them encompassed under the term \ac{FDA}, which essentially explores statistical techniques to process, model and make inference on data varying over a continuum. A partial survey on such methods is \citet{cuevas2014partial}, while a more detailed exposition of the theory and applications can be found for example in \citet{hsing2015theoretical} and \citet{horvath2012inference}.

\acs{FDA} has emerged as one of the most active areas of statistical research over the past decades, driven by the increasing availability of data collected over time, space, or other continuous domains. In contrast to classical multivariate statistics, where each data point is a vector, the \acs{FDA} framework treats each observation as a curve, surface, or more general function, typically viewed as a realization of an underlying stochastic process. This paradigm shift has enabled more nuanced analyses in diverse fields including climatology, biomechanics, neuroimaging, and econometrics, where the fundamental units of observation are inherently functional~\citep{ullah2013applications}. The monograph by \citet{ramsay2005functional} is generally regarded as a standard reference for the subject, while detailed up-to-date accounts of the theory and applications can be found for example in \citet{hsing2015theoretical, horvath2012inference,wang2016functional}.

\acs{FDA} is undoubtedly an active area of research, which finds applications in a wide variety of fields, such as biomedicine, finance, meteorology or chemistry~\citep{ullah2013applications}. Accordingly, there are many recent contributions on how to tackle functional data problems, both from a theoretical and practical standpoint. Chief among them is the approach of reducing the problem to a finite-dimensional one, for example using a truncated basis expansion or spline interpolation methods~\citep{muller2005generalized, aguilera2013comparative}. At the same time, much effort has been put into building a sound theoretical basis for \acs{FDA}, generalizing different frequentist concepts to the infinite-dimensional framework. Examples of this endeavor include the definition of centrality measures and depth-based notions for functional data~\citep{lopez2009concept}, functional analysis of variance tests~\citep{cuevas2004anova}, a functional Mahalanobis distance~\citep{galeano2015mahalanobis, berrendero2020mahalanobis}, or an extension of Fisher's discriminant analysis for random functions~\citep{shin2008extension}, among many others.

As the name suggests, \acs{FDA} techniques are heavily inspired by functional analysis tools and methods: Hilbert spaces, linear operators, orthonormal bases, and so on. Incidentally, a notion that also intersects with the classical theory of machine learning and pattern recognition, and that has attained popularity in recent years, is that of \ac{RKHSs} and their applications in functional data problems~\citep{kupresanin2010rkhs, yuan2010reproducing, berrendero2018use}. On the other hand, Bayesian inference methods are ubiquitous in the realm of statistics, and though they also make use of random functions, the approach is slightly different from the \acs{FDA} case. While there are recent works that offer a Bayesian treatment of functional data~\citep[e.g.][]{shi2011gaussian}, there is still no systematic approach to Bayesian methodologies within \acs{FDA}. It is precisely at this relatively unexplored intersection between \acs{FDA} and Bayesian methods that our work is aimed.



\section{Essentials on reproducing kernel Hilbert spaces}

The methodology proposed in this work relies heavily on the use of \acs{RKHSs}, so we will briefly outline the main characteristics of these spaces from a probabilistic point of view \citep[for a more detailed account, see][]{berlinet2004reproducing}. Let us denote by \(K(t, s)= \mathbb E[X(t)X(s)]\) the covariance function (the ``kernel'') of the centered process \(X\), and in what follows suppose that it is continuous. To construct the corresponding \acs{RKHS} \(\Hcal(K)\), we start by defining the functional space \(\Hcal_0(K)\) of all finite linear combinations of evaluations of \(K\), that is,
\begin{equation}\label{eq:h0}
  \Hcal_0(K) = \left\{ f \in L^2[0,1]: \ f(\cdot) = \sum_{j=1}^p \beta_j K(t_j, \cdot), \ p \in \N, \ \beta_j \in \R, \ t_j \in [0, 1] \right\}.
\end{equation}
This space is endowed with the inner product \(\dotprod{f}{g}_K = \sum_{j, k} \beta_j \gamma_k K(t_j, s_k)\), given that \(f(\cdot)=\sum_j \beta_j K(t_j, \cdot) \) and \(g(\cdot)=\sum_k \gamma_k K(s_k, \cdot)\). Then, \(\Hcal(K)\) is defined to be the completion of \(\Hcal_0(K)\) under the norm induced by the scalar product \(\dotprod{\cdot}{\cdot}_K\). As it turns out, functions in this space satisfy the so-called \textit{reproducing property}: \(\dotprod{K(t, \cdot)}{f}_K = f(t)\) for all \(f \in \Hcal(K)\) and \(t \in [0, 1]\). An important consequence of this identity is that \(\Hcal(K)\) is a space of genuine functions and not of equivalence classes, since the values of the functions at specific points are in fact relevant, unlike in \(L^2\)-spaces.

Now, a particularly useful approach in statistics is to regard \(\Hcal(K)\) as an isometric copy of a well-known space related to \(X\). Specifically, via \textit{Loève's isometry} \citep{loeve1948fonctions} one can establish a congruence \(\Psi_X\) between \(\Hcal(K)\) and the linear span of the process, \(\mathcal L(X)\), in the space of all random variables with finite second moment, \(L^2(\Omega)\) \citep[see Lemma 1.1 in][]{lukic2001stochastic}. This isometry is essentially the completion of the correspondence
\begin{equation*}
  \sum_{j=1}^p \beta_j X(t_j) \longleftrightarrow \sum_{j=1}^p \beta_j K(t_j, \cdot),
\end{equation*}
and can be formally defined, in terms of its inverse, as \(\Psi^{-1}_X(U)(t) = \E[U X(t)]\) for \(U \in \mathcal L(X)\).
Despite the close connection between the process \(X\) and the space \(\Hcal(K)\), special care must be taken when dealing with concrete realizations, since in general the trajectories of \(X\) do not belong to the corresponding \acs{RKHS} with probability one \citep[][Corollary~7.1]{lukic2001stochastic}. As a consequence, the expression \(\dotprod{x}{f}_K\) is ill-defined and lacks meaning when \(x\) is a realization of \(X\). However, following Parzen's approach in his seminal work \citep[][Theorem~4E]{parzen1961approach}, we can leverage Loève's isometry and identify \(\dotprod{X}{f}_K\) with the random variable in \(\mathcal L(X)\) associated with \(f\in \Hcal(K)\), so that \(\dotprod{x}{f}_K \) with \(x=X(\omega)\) is defined as the image \(\Psi_x(f) := \Psi_X(f)(\omega)\). This notation, viewed as a formal extension of the inner product, often proves to be useful and convenient.

%*****************************************
%*****************************************
%*****************************************
%*****************************************
%*****************************************
